\documentclass[11pt]{article}
\usepackage{amsmath,amssymb,amsthm}

\newtheorem{theorem}{Theorem}
\newtheorem{corollary}[theorem]{Corollary}
\newtheorem{lemma}[theorem]{Lemma}

\newcommand{\Up}{U_+}
\newcommand{\rev}{\operatorname{rev}}

\title{Same-depth inverse rigidity across two preceding record atoms}
\author{}
\date{27 August 2026}

\begin{document}
\maketitle

For a positive odd integer $x$, put
\[
 \Up(x):=\frac{3x+1}{2^{v_2(3x+1)}}.
 \tag{1}
\]
For an exponent word $w=(a_1,\ldots,a_L)$, write
\[
 B(w):=\sum_{i=1}^L a_i,
 \qquad
 C(w):=\sum_{i=1}^L3^{L-i}2^{a_1+\cdots+a_{i-1}}.
 \tag{2}
\]
Thus
\[
 2^{B(w)}\Up^L(x)=3^Lx+C(w).
 \tag{3}
\]
On an interior deficit-four record tail, use the atom notation
\[
 A:=1,\qquad B:=12,\qquad C:=1113,\qquad D:=1122.
 \tag{4}
\]
The length and excess mass of a nontrivial atom are
\[
 \begin{array}{c|ccc}
 P&B&C&D\\ \hline
 |P|&2&4&4\\
 B(P)-|P|&1&2&2.
 \end{array}
 \tag{5}
\]

If $t$ is an odd target and $M\geq L$, let $\mathcal R_L(t;M)$
denote the reverse exponent strings $(d_1,\ldots,d_L)$ of total mass at
most $M$ which give a positive odd $L$-step predecessor of $t$.  They are
generated by
\[
 t_0=t,
 \qquad
 t_i=\frac{2^{d_i}t_{i-1}-1}{3},
 \tag{6}
\]
with every $t_i$ a positive odd integer.  For a recursion with $\ell$
steps still unprocessed, only $t_i\bmod 3^\ell$ is needed.

Fix $P,Q\in\{B,C,D\}$ and $r,s\geq0$, and consider the exact word
\[
 W_{P,Q}(r,s):=P A^r Q A^s C.
 \tag{7}
\]
Put
\[
 \begin{aligned}
 L&:=|W_{P,Q}(r,s)|,
 &B_0&:=B(W_{P,Q}(r,s)),\\
 \delta_P&:=\begin{cases}1,&P=C,\\0,&P\in\{B,D\},\end{cases}
 &M&:=B_0-\delta_P.
 \end{aligned}
 \tag{8}
\]

\begin{lemma}[finite reverse-profile certificate]
Let $q$ be a positive odd start for the exact word (7), and let
$y=\Up^L(q)$.  Then
\[
 \boxed{
 \mathcal R_L(y;M)=
 \begin{cases}
  \{\rev W_{P,Q}(r,s)\},&P\in\{B,D\},\\[2mm]
  \{\rev(1121\,A^rQ A^sC)\},&P=C.
 \end{cases}}
 \tag{9}
\]
In particular, the minimum inverse mass is $B_0$ for $P=B,D$ and
$B_0-1$ for $P=C$, and the minimum-mass reverse word is unique.
\end{lemma}

\begin{proof}
The reverse sets satisfy the finite recursion
\[
 \mathcal R_{\ell+1}(t;N)
 =\bigcup_{\substack{d\geq1,\ d\leq N-\ell\\
                     3\mid2^dt-1}}
 \left\{(d,\gamma):
 \gamma\in\mathcal R_\ell
 \left(\frac{2^dt-1}{3};N-d\right)\right\}.
 \tag{10}
\]
The excess $M-L$ is at most $6$.  Hence every exponent tested in (10)
is at most $7$.

Process the last $L-|P|$ reverse steps and retain only branches leaving
at least one unit of mass per step of $P$.  For each retained branch,
record
\[
 (u,h):=
 \left(t_{L-|P|}\bmod3^{|P|},
       M-\sum_{i=1}^{L-|P|}d_i\right).
 \tag{11}
\]
\newpage
The complete profiles are as follows; every displayed pair has
multiplicity one.
\[
\begin{array}{c|c}
(P,Q)&\text{profile }\Pi_{P,Q}(r,s)\\ \hline
(B,B)&\{(4,3),(0,2)\}\\
(B,C)&\{(6,4),(7,2),(4,3)\}\\
(B,D)&\{(4,3),(0,2)\}\\[1mm]
(C,B)&\{(20,5),(41,4)\}\\
(C,C)&\{(50,6),(39,4),(20,5)\}\\
(C,D)&\{(20,5),(41,4)\}\\[1mm]
(D,B)&\{(10,6),(41,4),(21,5)\}\\
(D,C)&\{(45,7),(19,5),(10,6),(41,4)\}\\
(D,D)&\{(43,4),(10,6),(41,4),(21,5)\}.
\end{array}
\tag{12}
\]
When $(P,Q)=(D,C)$ and $r=0$, there is one additional pair $(39,4)$.

Here is an exact induction certificate for (12).  At $r=s=0$, the
triples $(L,M,y\bmod3^L)$ are
\[
\begin{array}{c|ccc}
PQ&BB&BC&BD\\ \hline
(L,M,y_0)&(8,12,2045)&(10,15,33959)&(10,15,51374)
\end{array}
\tag{13}
\]
\[
\begin{array}{c|ccc}
PQ&CB&CC&CD\\ \hline
(L,M,y_0)&(10,14,10064)&(12,17,375131)&(12,17,274448)
\end{array}
\tag{14}
\]
\[
\begin{array}{c|ccc}
PQ&DB&DC&DD\\ \hline
(L,M,y_0)&(10,15,37037)&(12,18,73325)&(12,18,504083).
\end{array}
\tag{15}
\]
Starting from these residues, substitution of $1\leq d\leq7$ in
(10) gives precisely the nine rows of (12).  Inserting one additional
letter $A$ in either run raises both $L$ and $M$ by one.  More
explicitly, if $w=UV$, $b=B(U)$, and $k=|V|$, then
\[
 C(UAV)=3^{k+1}C(U)+2^b\bigl(3^k+2C(V)\bigr).
\tag{16}
\]
Substitution of this identity in (10) leaves every profile in (12)
unchanged under $r\mapsto r+1$ and $s\mapsto s+1$, with one exception.
Writing $\Pi^*_{D,C}$ for the four-pair set displayed in the $(D,C)$
row, that exception is
\[
 \begin{aligned}
 \Pi_{D,C}(0,s)&=\Pi^*_{D,C}\cup\{(39,4)\},\\
 \Pi_{D,C}(r,s)&=\Pi^*_{D,C}\qquad(r\geq1).
 \end{aligned}
 \tag{17}
\]
The $s$-lift preserves both alternatives in (17), while the first
$r$-lift deletes $(39,4)$.  Equations (13)--(17) are a two-variable
induction proving (12).  This
certificate involves only division by $3$, reduction modulo $9$ or
$81$ at the terminal boundary, and the seven possible values of $d$;
there is no unlisted asymptotic or positivity step.

It remains to complete the first atom.  Direct use of (10) gives
\[
\begin{array}{c|c}
P&\text{nonempty completion among all pairs in (12)}\\ \hline
B&\mathcal R_2(4;3)=\{(2,1)\},\\
C&\mathcal R_4(20;5)=\{(1,2,1,1)\},\\
D&\mathcal R_4(10;6)=\{(2,2,1,1)\}.
\end{array}
\tag{18}
\]
All other completions are empty; explicitly,
\[
\begin{aligned}
 &\mathcal R_2(0;2)=\mathcal R_2(6;4)
  =\mathcal R_2(7;2)=\varnothing,\\
 &\mathcal R_4(41;4)=\mathcal R_4(50;6)
  =\mathcal R_4(39;4)=\varnothing,\\
 &\mathcal R_4(21;5)=\mathcal R_4(45;7)
  =\mathcal R_4(19;5)=\mathcal R_4(43;4)=\varnothing.
\end{aligned}
\tag{19}
\]
In (18)--(19), a residue denotes any odd target in that residue class
modulo $3^{|P|}$; recursion (10) depends on no further information.

The unique retained suffix leading to $(4,3)$ or $(10,6)$ is the reverse
of $A^rQA^sC$, and (18) appends the reverse of $B$ or $D$.  The unique
suffix leading to $(20,5)$ is again the actual suffix, while its
completion $(1,2,1,1)$ is the reverse of $1121$.  This proves (9).
\end{proof}

\begin{theorem}[two-preceding-atom same-depth minimum]
Let $q$ be the start and $y$ the endpoint of the exact word
$W_{P,Q}(r,s)$ in (7).  Among all positive odd $L$-step predecessors of
$y$, the unique least one is
\[
 \boxed{
 q_{\min}=
 \begin{cases}
 q,&P\in\{B,D\},\\[1mm]
 (q-1)/2,&P=C.
 \end{cases}}
 \tag{20}
\]
When $P=C$, the second line is exactly the already known half-coalescing
seed at the first atom $C$; its word is $1121A^rQA^sC$.
\end{theorem}

\begin{proof}
Lemma~1 gives the minimum inverse mass $M$ and its unique word.  For
$P=C$, the exact local identity
\[
 \Up^4\!\left(\frac{q-1}{2}\right)=\Up^4(q)
 \tag{21}
\]
replaces the first word $1113$ by $1121$, proving that the start in
(20) is attached to that minimum word.

We first record coarse lower bounds on this start.  Exactness of the
first atom gives
\[
 \begin{aligned}
 q&\equiv11\pmod {16}&& (P=B),\\
 q&\equiv79\pmod {128}&& (P=C),\\
 q&\equiv71\pmod {128}&& (P=D).
 \end{aligned}
 \tag{22}
\]
Every suffix $A^rQA^sC$ starts with exponent $1$.  Direct iteration of
the finitely many small representatives gives
\[
\begin{array}{c|c|c}
P&q&\text{initial exponent or atom pattern}\\ \hline
B&11&123\\
 &27&BA^2DB\\
 &43&122\\
 &59&1214\\
C&79&11134\\
D&71&DBAB.
\end{array}
\tag{23}
\]
None has the form $PA^rQA^sC$.  Therefore
\[
 q>64\quad(P=B),
 \qquad
 \frac{q-1}{2}>64\quad(P=C),
 \qquad
 q>128\quad(P=D).
 \tag{24}
\]

Let $q_*$ denote the value in (20), and let $C_*$ be the carry of its
minimum word.  From (5) and (8), the bounds in (24) imply
\[
 \boxed{q_*>2^{M+1-L}.}
 \tag{25}
\]
Indeed, the right side is at most $64$ for $P=B,C$ and at most $128$
for $P=D$.

Now let $v=(c_1,\ldots,c_L)$ be any other exact word from a positive odd
$z$ to $y$, and put $B'=B(v)$.  Lemma~1 gives $B'\geq M+1$.  Positivity
of the exponents gives the uniform carry estimate
\[
 C(v)\leq2^{B'-L}(3^L-2^L)<2^{B'-L}3^L.
 \tag{26}
\]
The minimum affine equation and (25) imply $y>(3/2)^L$.  Hence
\[
\begin{aligned}
 z
 &=\frac{2^{B'}y-C(v)}{3^L}\\
 &>\frac{2^{B'}y}{3^L}-2^{B'-L}\\
 &\geq\frac{2^{M+1}y}{3^L}-2^{M+1-L}\\
 &=2q_*+\frac{2C_*}{3^L}-2^{M+1-L}>q_*.
\end{aligned}
\tag{27}
\]
Thus the minimum-mass start is also the unique least positive start.
\end{proof}

\begin{corollary}[two preceding atoms give no new same-depth descent]
Suppose an interior record tail contains
\[
 P A^r Q A^s C,
 \qquad P,Q\in\{B,C,D\}.
 \tag{28}
\]
If $P=B$ or $D$, no same-depth inverse branch starts below the actual
block start.  If $P=C$, the unique least branch is only the local
half-seed $(q-1)/2$ belonging to that earlier $C$ atom.  Consequently a
new pullback of the terminal $C$ half-coalescence cannot produce a
same-depth descent while crossing only its two preceding nontrivial
atoms.  Such a construction must change raw inverse depth or cross at
least three preceding nontrivial atoms.
\end{corollary}

\begin{proof}
Apply Theorem~2.  Its $P=C$ alternative is (20), so it introduces no
inverse mechanism beyond the already classified local $C$ half-seed.
\end{proof}

\end{document}
